\section{Research guide: detection, evasion and
adaptation}\label{record-docs-research-guide-md-research-guide-detection-evasion-and-adaptation}

Commands in this guide run from the repository root. Validation records
describe their original runs; see the \href{../README.md}{project
overview} for the reading path.

This project is a research workbench for three related questions:

\begin{enumerate}
\def\labelenumi{\arabic{enumi}.}
\tightlist
\item
  Can a model distinguish AI-generated content from human-created
  content?
\item
  Can a generator learn to evade that detector while preserving the
  requested meaning or image prompt?
\item
  What happens when the detector is retrained on the attacker's new
  outputs?
\end{enumerate}

It supports text and image detectors, detector ensembles, text and image
evasion methods, repeated attacker and defender updates, leakage-aware
benchmarks, and a Streamlit scoring application.

AI-content detection is statistical evidence, not proof of authorship. A
score can change with language, topic, generator, decoding settings,
image processing, and adversarial rewriting. Use this project to run
controlled experiments or to assist review, not as the sole basis for a
consequential decision.

The bibliography is in
\href{../ai_content_detector/references.bib}{\texttt{\seqsplit{references.bib}}}.
This README explains the experimental workflow; the
\href{records/DETECTORS.tex}{detector catalogue} explains individual
detection methods without requiring the reader to first read those
papers.

\subsection{The problem in plain
language}\label{record-docs-research-guide-md-the-problem-in-plain-language}

AI generators leave patterns, but no single pattern is universal. A text
model might produce unusually predictable token sequences. An image
generator might leave frequency artifacts or react differently to
denoising. A supervised classifier can learn such patterns from
examples, while a zero-shot detector can derive a statistic from a
pretrained model without training a new classifier.

An adaptive attacker makes the problem harder. It observes a detector
score and changes the generator to lower that score. A defender can then
retrain on these new evasive samples. This repository represents that
loop as:

\begin{Verbatim}
prompt -> generator -> candidate content -> detector -> AI score
             ^                              |
             |________ evasion reward ______|

fixed evaluation set -> measure attacker -> retrain defender -> measure again
\end{Verbatim}

The key experimental question is therefore not only ``does the detector
work on today's samples?'' It is also ``does it continue to work on
samples optimized against it, including samples from prompts and
detectors that training never saw?''

\subsection{Watermarking
experiments}\label{record-docs-research-guide-md-watermarking-experiments}

The \href{records/docs/instruct_watermarking.tex}{current text protocol}
tests TextSeal 2 and the stronger DistSeal-inspired latent student on
Qwen/Qwen3.5-4B instruct. BIRA uses Llama-3.1-8B-Instruct. Direct
probability resampling and confidence smoothing are also tested, with
ordinary paraphrasing as a control. Quality is measured separately from
raw evasion. Earlier Base-model watermark experiments were removed;
TextSeal 1 is retired. Run
\texttt{\seqsplit{scripts/instruct\_watermark\_experiments.py}} for the
replacement experiment. Frozen snapshots, checkpoint state and Slurm
advance signals support resumption on two-hour dev QoS.

The \href{records/docs/watermarking.tex}{preserved image protocol} uses
released RAR-XL and UViT-H/DC-AE DistSeal checkpoints. Its removal jobs
test clean-checkpoint merging and clean-data LoRA/full training, with
separate worker contracts and owner scoring. Native image fitting uses
800 classes and reports separately on the 200 classes excluded from
fitting. These image artifacts were preserved.

\subsection{Essential
vocabulary}\label{record-docs-research-guide-md-essential-vocabulary}

{\scriptsize\setlength{\tabcolsep}{3pt}\def\LTcaptype{none} % do not increment counter
\begin{longtable}[]{@{}
  >{\raggedright\arraybackslash}p{(\linewidth - 2\tabcolsep) * \real{0.5000}}
  >{\raggedright\arraybackslash}p{(\linewidth - 2\tabcolsep) * \real{0.5000}}@{}}
\toprule\noalign{}
\begin{minipage}[b]{\linewidth}\raggedright
Term
\end{minipage} & \begin{minipage}[b]{\linewidth}\raggedright
Meaning in this project
\end{minipage} \\
\midrule\noalign{}
\endhead
\bottomrule\noalign{}
\endlastfoot
\textbf{Detector} & A component that maps text or an image to an
AI-oriented score and a label. \\
\textbf{Supervised detector} & A classifier trained on labeled human and
AI examples. It needs a fitted checkpoint. \\
\textbf{Zero-shot detector} & A method that computes a statistic using
pretrained models without fitting a task-specific detector classifier.
It can still require large model weights and a threshold. \\
\textbf{Raw statistic} & The detector's native measurement, such as a
likelihood ratio, curvature, cosine similarity, or anomaly score. \\
\textbf{Calibration} & Learning a mapping from scores to probabilities
on a separate representative validation set. \\
\textbf{Evasion attack} & A rewrite or generator update intended to
reduce detection while retaining useful content. \\
\textbf{Attacker} & The generator or transformation being optimized. \\
\textbf{Defender} & The detector being evaluated or retrained. \\
\textbf{Held-out} & Data or a detector excluded from training and used
only to test generalization. \\
\textbf{Checkpoint} & Saved trained weights or a fitted classical model.
Large project checkpoints are not stored in git. \\
\end{longtable}
}

\subsection{Quickstart}\label{record-docs-research-guide-md-quickstart}

Run every command from the repository root. The project uses the root
\texttt{\seqsplit{uv}} environment and expects
\texttt{\seqsplit{PYTHONPATH=.}} for module imports.

\subsubsection{1. Create the
environment}\label{record-docs-research-guide-md-create-the-environment}

\begin{Verbatim}
uv venv
source .venv/bin/activate
uv pip install -r requirements.txt
uv run python -m spacy download en_core_web_sm
\end{Verbatim}

On the CUDA 12.4 research VM, install the GPU environment instead:

\begin{Verbatim}
uv pip install -r requirements_gpu.in
\end{Verbatim}

For Qwen3.5 RL on the Jean Zay H100 environment, optional
linear-attention and causal-convolution kernels avoid the slow PyTorch
fallbacks. Keep the installed PyTorch/CUDA versions and build the
convolution extension against that stack:

\begin{Verbatim}
module load cuda/12.4.1
export CC=/usr/bin/gcc CXX=/usr/bin/g++ CUDAHOSTCXX=/usr/bin/g++
uv pip install ninja wheel
CAUSAL_CONV1D_FORCE_BUILD=TRUE MAX_JOBS=4 TORCH_CUDA_ARCH_LIST=9.0 \
  uv pip install --no-build-isolation --no-deps -r requirements_rl_kernels.in
\end{Verbatim}

Run GPU validation inside Slurm. The official prebuilt convolution
wheels failed the ABI import check with this PyTorch build. The kernel
pins and native build preserve PyTorch 2.6/CUDA 12.4 rather than
upgrading the experiment environment.

The standard installation is enough for the CPU and mock-based test
suite. Most large detectors and all meaningful training runs need
downloaded Hugging Face weights, project checkpoints, and usually a GPU.
On macOS, XGBoost also requires an OpenMP runtime such as
\texttt{\seqsplit{libomp}}.

\subsubsection{2. Check imports and
wiring}\label{record-docs-research-guide-md-check-imports-and-wiring}

\begin{Verbatim}
PYTHONPATH=. uv run python -m ai_content_detector.rl_evasion.run_experiments \
  --experiment smoke_test
\end{Verbatim}

This is deliberately lightweight. It checks imports, configuration
objects, reward plumbing, and trainer construction. It does not download
or run every large model.

\subsubsection{3. Run the focused test
suite}\label{record-docs-research-guide-md-run-the-focused-test-suite}

\begin{Verbatim}
PYTHONPATH=. uv run python -m pytest ai_content_detector/tests -q
\end{Verbatim}

These tests validate formulas, tensor shapes, label conventions, split
contracts, failure behavior, seeding, and leakage controls using
CPU-scale fixtures and mocks.

\subsubsection{4. Launch the scoring
app}\label{record-docs-research-guide-md-launch-the-scoring-app}

\begin{Verbatim}
PYTHONPATH=. uv run streamlit run ai_content_detector/app.py
\end{Verbatim}

Open the URL printed by Streamlit, select the text or image tab, choose
the available detectors, provide content, and run the analysis. The app
reports why a detector is unavailable. It does not replace a missing
checkpoint with a heuristic result. At least one usable detector is
required.

The app treats text shorter than 50 characters as too short for reliable
use. Some individual methods impose stronger requirements. For example,
Disrupt-and-Recover defaults to at least 35 words.

\subsubsection{5. Run a small
benchmark}\label{record-docs-research-guide-md-run-a-small-benchmark}

Once Hugging Face downloads are available, the following command
downloads 20 HC3 examples and runs CPU-compatible attack baselines
against a RoBERTa detector. It is a practical pipeline check, not a
publishable benchmark:

\begin{Verbatim}
PYTHONPATH=. uv run python \
  -m ai_content_detector.rl_evasion.benchmarking.benchmark \
  --dataset hc3 \
  --max-samples 20 \
  --lightweight \
  --detectors roberta_classifier \
  --output-dir results/benchmark_quickstart
\end{Verbatim}

\texttt{\seqsplit{--lightweight}} changes the attack methods, not the
scoring stack. This command still downloads the RoBERTa detector, the E5
semantic encoder, HC3, and NLTK resources. Use the smoke test and unit
tests for an offline CPU check.

HC3 contains aligned prompts, human answers, and ChatGPT answers. The
CNN/DailyMail option contains prompts and human continuations but no
stored AI continuations, so it can evaluate generation methods but skips
rewrite-only methods that require an existing AI text.

\subsection{Repository
map}\label{record-docs-research-guide-md-repository-map}

\begin{Verbatim}
ai_content_detector/
[U+251C][U+2500][U+2500] app.py                              # Streamlit text and image UI
[U+251C][U+2500][U+2500] detectors/
[U+2502]   [U+251C][U+2500][U+2500] ensemble.py                    # Detector interface and aggregation
[U+2502]   [U+251C][U+2500][U+2500] text_detectors.py              # Text detectors
[U+2502]   [U+251C][U+2500][U+2500] image_detectors.py             # Image detectors
[U+2502]   [U+2514][U+2500][U+2500] style_detector.py              # LUAR few-shot style detector
[U+251C][U+2500][U+2500] rl_evasion/
[U+2502]   [U+251C][U+2500][U+2500] config.py                      # Reproducible experiment configs
[U+2502]   [U+251C][U+2500][U+2500] run_experiments.py             # Unified experiment CLI
[U+2502]   [U+251C][U+2500][U+2500] text_evasion/                  # GRPO, SPIN, rewards, features
[U+2502]   [U+251C][U+2500][U+2500] image_evasion/                 # DDPO training and evaluation
[U+2502]   [U+251C][U+2500][U+2500] arms_race/                     # Repeated adaptation and MAML
[U+2502]   [U+2514][U+2500][U+2500] benchmarking/                  # Datasets, baselines, runner
[U+251C][U+2500][U+2500] references.bib
[U+2514][U+2500][U+2500] tests/
\end{Verbatim}

Internal supervised detectors reuse checkpoint loading and model
wrappers from \texttt{\seqsplit{ml\_pipeline/}}. Shared checkpoint
behavior should remain there rather than being duplicated in this
subproject.

\subsection{Detector interface and score
semantics}\label{record-docs-research-guide-md-detector-interface-and-score-semantics}

Every detector implements \texttt{\seqsplit{BaseDetector}}:

\begin{Verbatim}
result = detector.detect(content)

result.score          # float in [0, 1], oriented so larger means more AI-like
result.label          # "human", "ai", "uncertain", or "error" in an ensemble
result.detector_name  # runtime detector name
result.details        # raw statistic, thresholds, semantics, or failure context
\end{Verbatim}

\texttt{\seqsplit{is\_available()}} checks dependencies, model
configuration, support data, and checkpoints.
\texttt{\seqsplit{unload()}} releases known model references and clears
CUDA cache when possible. Heavy models are loaded lazily.

\subsubsection{A score is not automatically a
probability}\label{record-docs-research-guide-md-a-score-is-not-automatically-a-probability}

All public scores use the same direction, with \texttt{\seqsplit{0}}
more human-like and \texttt{\seqsplit{1}} more AI-like. That common
range makes display and aggregation convenient, but it does not give
every number the same interpretation.

\begin{itemize}
\tightlist
\item
  \texttt{\seqsplit{classifier\_probability}} comes from a fitted
  classifier's softmax or \texttt{\seqsplit{predict\_proba}}. It is
  still only as reliable as that classifier's data.
\item
  \texttt{\seqsplit{calibrated\_probability}} means a separate
  calibrator was fitted on labeled validation examples.
\item
  \texttt{\seqsplit{uncalibrated\_*}} means the value is a normalized or
  monotonic display score. It preserves ordering but must not be read as
  confidence.
\item
  A raw statistic and its decision threshold are retained in
  \texttt{\seqsplit{details}} when a paper method makes its decision in
  the raw space.
\end{itemize}

For illustration, suppose Binoculars computes a raw ratio of
\texttt{\seqsplit{0.77}}. Its display transform is
\texttt{\seqsplit{sigmoid((0.9015 - raw) / 0.05)}}, which maps
\texttt{\seqsplit{0.77}} to \texttt{\seqsplit{0.93}} and the paper
threshold itself to exactly \texttt{\seqsplit{0.5}}, and
\texttt{\seqsplit{details["score\_semantics"]}} is
\texttt{\seqsplit{uncalibrated\_monotonic\_transform}}. Nothing was
fitted; the map was chosen to put \texttt{\seqsplit{0.5}} at the
threshold so the display score stays readable around the decision point.
Read \texttt{\seqsplit{0.93}} as ``further below the threshold than
\texttt{\seqsplit{0.60}} would be'', not as ``93 percent probability''.

Inspect \texttt{\seqsplit{result.details["score\_semantics"]}} whenever
it is present.

The generic three-way helper labels scores at least
\texttt{\seqsplit{0.65}} as \texttt{\seqsplit{ai}}, at most
\texttt{\seqsplit{0.35}} as \texttt{\seqsplit{human}}, and the interval
between them as \texttt{\seqsplit{uncertain}}. Methods with a native
paper threshold make their label from the raw paper statistic instead,
which is why Binoculars and Fast-DetectGPT never return
\texttt{\seqsplit{uncertain}}.

\subsubsection{Ensemble
behavior}\label{record-docs-research-guide-md-ensemble-behavior}

\texttt{\seqsplit{EnsembleAggregator}} runs all selected detectors
independently. Its default weighted average is

\[
S_{\mathrm{ensemble}} = \sum_{j=1}^{m} \tilde{w}_j S_j,
\qquad
\tilde{w}_j = \frac{w_j}{\sum_k w_k}.
\]

For illustration, suppose three equally weighted detectors return the
following outputs. These are hypothetical scores, not measured results:

{\scriptsize\setlength{\tabcolsep}{3pt}\def\LTcaptype{none} % do not increment counter
\begin{longtable}[]{@{}
  >{\raggedright\arraybackslash}p{(\linewidth - 4\tabcolsep) * \real{0.3333}}
  >{\raggedright\arraybackslash}p{(\linewidth - 4\tabcolsep) * \real{0.3333}}
  >{\raggedright\arraybackslash}p{(\linewidth - 4\tabcolsep) * \real{0.3333}}@{}}
\toprule\noalign{}
\begin{minipage}[b]{\linewidth}\raggedright
detector
\end{minipage} & \begin{minipage}[b]{\linewidth}\raggedright
score
\end{minipage} & \begin{minipage}[b]{\linewidth}\raggedright
label
\end{minipage} \\
\midrule\noalign{}
\endhead
\bottomrule\noalign{}
\endlastfoot
Binoculars & 0.067 & human \\
TF-IDF + logistic regression & 0.62 & uncertain \\
ModernBERT & 0.83 & ai \\
\end{longtable}
}

The average is \texttt{\seqsplit{(0.067 + 0.62 + 0.83) / 3 = 0.506}},
but the aggregate \textbf{label} comes from votes, not from that number:
one \texttt{\seqsplit{ai}} vote and one \texttt{\seqsplit{human}} vote
cancel, so the label is \texttt{\seqsplit{uncertain}} even though the
score sits above \texttt{\seqsplit{0.5}}. The two rules answer different
questions, which is why \texttt{\seqsplit{aggregate\_score}} and
\texttt{\seqsplit{aggregate\_label}} are reported separately.

Without a fitted logistic calibrator, this is an uncalibrated normalized
score and the label comes from detector votes. With a calibrator fitted
on disjoint labeled validation data, the aggregate is reported as a
calibrated probability and the label comes from a
\texttt{\seqsplit{0.5}} cut. An empty ensemble, a calibration set with
one class, and a run in which every detector fails all raise clear
errors. A failure from one detector appears in the per-detector results
with label \texttt{\seqsplit{error}} and is excluded from the aggregate.

\subsection{Detector
catalogue}\label{record-docs-research-guide-md-detector-catalogue}

See the \href{records/DETECTORS.tex}{text and image detector catalogue}
for individual methods, worked illustrations, formulas, and
implementation requirements. The sections below focus on evasion,
adaptation, and evaluation.

\subsection{Evasion
methods}\label{record-docs-research-guide-md-evasion-methods}

\subsubsection{The reward
design}\label{record-docs-research-guide-md-the-reward-design}

Text and image attacks optimize more than detector evasion. If the only
reward were \texttt{\seqsplit{1 - detector\_score}}, a generator could
emit empty or irrelevant content. The text configuration therefore
combines detector evasion, semantic similarity, and a length and quality
term:

\[
R_{\mathrm{text}} = 0.6R_{\mathrm{evasion}}
                  + 0.3R_{\mathrm{semantic}}
                  + 0.1R_{\mathrm{quality}}.
\]

The image reward is
\texttt{\seqsplit{0.64 * evasion + 0.26 * CLIP alignment + 0.10 * reference}}.
Reference similarity is \texttt{\seqsplit{1 / (1 + mean\_KL / scale)}},
where conditional KL compares stochastic DDIM transitions with initial
SD1.5, averaged over latent coordinates and stochastic timesteps for
each image. All three components are bounded in
\texttt{\seqsplit{[0, 1]}} before joint advantage normalization; there
is no extra direct KL loss. The dedicated experiment freezes the scale
from a training-only probe and then restarts the main attacker from
initial SD1.5. See the
\href{records/docs/image_reference_reward.tex}{reference reward
protocol} for calibration, observed-scale diagnostics and the
distinction from exact KL regularization.

\subsubsection{One example used by the three text
methods}\label{record-docs-research-guide-md-one-example-used-by-the-three-text-methods}

The text methods below use the same illustrative item. Generated
completions, detector scores, similarities, and rewards are hypothetical
examples, not measured results. Dataset loading and prompt construction
are described in the implementation; this question-and-answer item is
for explanation only:

\begin{itemize}
\tightlist
\item
  \textbf{Prompt:}
  \texttt{\seqsplit{Explain why sourdough bread rises.}}
\item
  \textbf{Human reference} (illustrative):
  \texttt{\seqsplit{Wild yeast and bacteria eat the sugars in the flour and give off carbon dioxide, and the gluten network traps that gas. In the oven the starch sets, so the loaf keeps the shape the gas gave it.}}
\item
  \textbf{Detector:} any detector that returns an AI-oriented score in
  \([0, 1]\), for example \texttt{\seqsplit{BinocularsDetector}}. Lower
  means less AI-like according to that detector.
\end{itemize}

\subsubsection{GRPO text
evasion}\label{record-docs-research-guide-md-grpo-text-evasion}

The model learns from rewards on its own generated continuations. This
implementation still uses training prompts and human reference
continuations from a corpus: the references supply semantic-similarity
and length rewards. It does not require pre-labelled preference pairs.

\texttt{\seqsplit{GRPOTextEvasionTrainer}} samples
\texttt{\seqsplit{grpo\_num\_generations = 4}} completions for the
prompt and scores each with the reward above (including the length
term):

{\scriptsize\setlength{\tabcolsep}{3pt}\def\LTcaptype{none} % do not increment counter
\begin{longtable}[]{@{}
  >{\raggedright\arraybackslash}p{(\linewidth - 6\tabcolsep) * \real{0.2500}}
  >{\raggedright\arraybackslash}p{(\linewidth - 6\tabcolsep) * \real{0.2500}}
  >{\raggedright\arraybackslash}p{(\linewidth - 6\tabcolsep) * \real{0.2500}}
  >{\raggedright\arraybackslash}p{(\linewidth - 6\tabcolsep) * \real{0.2500}}@{}}
\toprule\noalign{}
\begin{minipage}[b]{\linewidth}\raggedright
completion
\end{minipage} & \begin{minipage}[b]{\linewidth}\raggedright
detector score
\end{minipage} & \begin{minipage}[b]{\linewidth}\raggedright
semantic
\end{minipage} & \begin{minipage}[b]{\linewidth}\raggedright
reward
\end{minipage} \\
\midrule\noalign{}
\endhead
\bottomrule\noalign{}
\endlastfoot
``Sourdough rises because wild yeast ferments the sugars in flour and
produces carbon dioxide gas, which expands inside the dough during
baking.'' & 0.81 & 0.94 & 0.48 \\
``Bakers mix water and starter into the flour, and tiny organisms
consume the sugars and exhale gas that inflates the dough.'' & 0.44 &
0.91 & 0.69 \\
``It rises.'' & not scored & not scored & 0.00 (too short) \\
``The leavening cultures feed on the flour, the dough traps the gas they
emit, and heat fixes the risen structure.'' & 0.52 & 0.93 & 0.64 \\
\end{longtable}
}

The reward is normalized \textbf{inside the group}, so each completion
is compared with the group mean: \(A_i = (r_i - \bar{r}) / \sigma_r\).
The best completion gets a positive advantage and the worst a negative
one, encouraging higher-reward continuations. Outputs shorter than 20
whitespace-separated words receive zero reward before detector or
semantic scoring. Semantic and length rewards further discourage
irrelevant or degenerate outputs, without guaranteeing quality.

No critic and no value network: the group average replaces them. The
native TRL arguments are saved in
\texttt{\seqsplit{trl\_training\_arguments.json}}; its existing default
KL coefficient is \texttt{\seqsplit{beta=0}}. The dataclass's
\texttt{\seqsplit{grpo\_kl\_coeff=0.05}} applies only to the disabled
manual fallback. The generator is adapted with LoRA. The method is
called GRPO only when the TRL implementation is available; the local
fallback is explicitly named group-normalized REINFORCE and is disabled
by default.

\subsubsection{SPIN with feature
monitoring}\label{record-docs-research-guide-md-spin-with-feature-monitoring}

No reward and no detector: the model learns from its own samples plus
the human continuation.

\texttt{\seqsplit{SPINTrainerWithFeatureMonitoring}} alternates two
steps on the same item.

\begin{enumerate}
\def\labelenumi{\arabic{enumi}.}
\tightlist
\item
  \textbf{Generate.} The current model writes its own continuation,
  typically with a high temperature so it drifts:
  \texttt{\seqsplit{Sourdough rises due to the activity of yeast and bacteria present in the starter, which metabolize sugars into gas during fermentation, causing the dough to expand as it bakes.}}
\item
  \textbf{Prefer.} The human continuation is \texttt{\seqsplit{chosen}},
  the generated one is \texttt{\seqsplit{rejected}}, and the loss raises
  the human continuation's relative log-probability:
\end{enumerate}

\[
\mathcal{L}_{\mathrm{SPIN}} = -\mathbb{E}\left[\log \sigma\left(\beta\left(
  \left[\log \pi_\theta(y_h) - \log \pi_{\mathrm{ref}}(y_h)\right]
- \left[\log \pi_\theta(y_g) - \log \pi_{\mathrm{ref}}(y_g)\right]\right)\right)\right].
\]

With \(\beta = 0.1\) (\texttt{\seqsplit{beta\_spin}}), the update simply
moves probability mass toward ``wild yeast and bacteria'' phrasing and
away from the flat textbook sentence, without any detector ever being
called during training. The frozen reference snapshot advances after
each SPIN iteration. Log-probabilities are conditional on the prompt and
cover the continuation only, so prompt tokens do not dominate the
signal.

Stylometric statistics, E5 embeddings, LUAR style embeddings, and
learned probe features are monitoring outputs. They are not gradient
losses because decoding text breaks a differentiable path from those
features back to model parameters. \texttt{\seqsplit{MultiSPINTrainer}}
remains as a compatibility alias.

\subsubsection{Adversarial
paraphrasing}\label{record-docs-research-guide-md-adversarial-paraphrasing}

No training at all: the detector is called at every generated token.

\texttt{\seqsplit{AdversarialParaphrasingBaseline}} rewrites the example
text token by token. At each position it takes the candidates from top-p
\texttt{\seqsplit{0.99}}, keeps at most top-k \texttt{\seqsplit{50}},
decodes the partial paraphrase for each one, and asks the detector to
score it.

An illustrative candidate comparison (assuming the guidance detector can
score these short prefixes):

{\scriptsize\setlength{\tabcolsep}{3pt}\def\LTcaptype{none} % do not increment counter
\begin{longtable}[]{@{}
  >{\raggedright\arraybackslash}p{(\linewidth - 4\tabcolsep) * \real{0.3333}}
  >{\raggedright\arraybackslash}p{(\linewidth - 4\tabcolsep) * \real{0.3333}}
  >{\raggedright\arraybackslash}p{(\linewidth - 4\tabcolsep) * \real{0.3333}}@{}}
\toprule\noalign{}
\begin{minipage}[b]{\linewidth}\raggedright
candidate token
\end{minipage} & \begin{minipage}[b]{\linewidth}\raggedright
partial paraphrase
\end{minipage} & \begin{minipage}[b]{\linewidth}\raggedright
detector score
\end{minipage} \\
\midrule\noalign{}
\endhead
\bottomrule\noalign{}
\endlastfoot
``Wild'' & \texttt{\seqsplit{Wild}} & 0.46 \\
``Naturally'' & \texttt{\seqsplit{Naturally}} & 0.38 \\
``The'' & \texttt{\seqsplit{The}} & 0.79 \\
``In'' & \texttt{\seqsplit{In}} & 0.71 \\
\end{longtable}
}

The lowest score wins, so the step commits to ``Naturally'', and the
search continues from that prefix. The detector is in the loop at every
token, which is why this is far more expensive than ordinary sampling
and why it requires a real guidance detector.

The benchmark also includes simpler controls: no transformation, neural
paraphrasing, WordNet synonym substitution, human-style prompt
instructions, high-temperature sampling, and optional OpenRouter
generation.

GRPO queries detectors during training; SPIN trains on human/generated
continuation pairs without detector queries. Both generate normally
afterwards. Adversarial paraphrasing pays at every token of every
generation.

\subsubsection{DDPO image
evasion}\label{record-docs-research-guide-md-ddpo-image-evasion}

\texttt{\seqsplit{DDPOImageEvasionTrainer}} fine-tunes a diffusion image
generator in three steps:

\begin{enumerate}
\def\labelenumi{\arabic{enumi}.}
\tightlist
\item
  \textbf{Generate} images from prompts, storing each image's sequence
  of denoising decisions (50 steps by default).
\item
  \textbf{Score} the final images and their trajectories: 64\% detector
  evasion, 26\% CLIP prompt alignment and 10\% bounded initial-model
  reference similarity.
\item
  \textbf{Update} the generator's UNet LoRA weights to make decisions
  associated with above-baseline rewards more likely, and those with
  below-baseline rewards less likely. DDPO-IS uses a clipped
  policy-gradient objective. The reference reward is computed during
  rollout with initial SD1.5 evaluated with LoRA disabled. Deterministic
  terminal transitions are excluded from policy updates and KL.
\end{enumerate}

Repeat this loop to favor images that evade the detectors while matching
the prompt. ``More human-like'' means lower detector AI scores here;
human-perceived realism and image quality are not guaranteed.

The DDPO default defender is now DDA. The original ArtiFact-trained
EfficientNet and CLIP classifiers remain available as baselines. The
unvalidated third-party SigLIP detector has been removed.
\texttt{\seqsplit{DDADetector}} loads the pinned official model and
source described by
\texttt{\seqsplit{results/research\_workflow/image\_dda/resources/provenance.json}}.
Its native preprocessing uses a 336-pixel center crop and CLIP
normalization. \texttt{\seqsplit{CommunityForensicsDetector}} is also
exported, with pinned official 224/384 checkpoints, native resizing and
ImageNet normalization. It loads lazily and is not included in the DDA
experiment, reward ensemble, or automatic downloads.

The removal passed 26 focused image detector, numerical, DDPO and
configuration tests. Streamlit startup also passed with detector
availability mocked as unavailable; model loading and scoring were not
exercised by that smoke check.

\subsubsection{Official DDA benchmark and one SD1.5
round}\label{record-docs-research-guide-md-official-dda-benchmark-and-one-sd1.5-round}

The dedicated driver evaluates official DDA on the existing ArtiFact
holdout, CIFAKE and Shoes without local detector training. It then warms
official DDA on 1,000 base-SD1.5 outputs paired with ArtiFact real
images. An internal 80/20 split reserves 200 pairs for validation; DDPO
uses only the 800 fitting prompts. After 50 DDPO epochs, one defender
update uses fresh attacked outputs on the same fitting pool. Both
defender fits use three epochs and learning rate 2e-5, updating only DDA
LoRA and classifier parameters and selecting validation loss.

The 500 evaluation prompts and 500 ArtiFact holdout real images remain
excluded from fitting. Before/after scoring reuses saved base and
attacked images. For the race, both classes are resized to 200 pixels
before the DDA center crop to control resolution differences. The
benchmark preserves native preprocessing. This is supervised adaptation
of released DDA, not a reproduction of its original aligned-data
training. Community Forensics is intentionally unused.

\begin{Verbatim}
PYTHONPATH=. uv --no-cache run --no-sync python scripts/image_experiments.py prepare
PYTHONPATH=. uv --no-cache run --no-sync python scripts/image_experiments.py submit MANIFEST
PYTHONPATH=. uv --no-cache run --no-sync python scripts/image_experiments.py report MANIFEST
\end{Verbatim}

The revised attacker first runs one training-only probe epoch without
reference reward, measures per-image KL after its updates, and fixes the
reference scale to the median positive measurement. Probe weights are
discarded. The main attacker still starts from initial SD1.5 and runs 50
epochs. Calibration, effective configuration and component variation
diagnostics are retained.

For an existing chain,
\texttt{\seqsplit{replace-attack PREVIOUS\_MANIFEST}} creates a new
frozen run that adopts its preparation, benchmark and warm-up jobs.
\texttt{\seqsplit{submit MANIFEST --held}} submits only replacement
attacker and defense jobs for dependency checks before swapping the old
pending jobs. It never cancels adopted jobs on a submission failure. The
\href{records/docs/image_reference_reward_results.tex}{replacement live
report} is separate from reports still written by the running original
prerequisites.

The five jobs use \texttt{\seqsplit{qos\_gpu\_h100-t3}}: an eight-hour
benchmark allocation and four-hour allocations for preparation, warm-up,
attack and defense. Dependencies prevent attack before warm-up and
defense before completed DDPO training. Frozen source copies, model
revisions, image identities, restart state and per-stage receipts are
retained. Preparation includes a real GPU forward/backward smoke check
before the production benchmark or race can start. The
\href{records/docs/image_reference_reward_results.tex}{live image
report} links the run's measured tables, CSV and attacker curves;
submitted work is not reported as a completed experiment.

Image workflow validation passed 51 focused tests covering native
preprocessing, AI-score direction, initial-reference KL and its
gradients, deterministic DDIM transitions, bounded rewards, probe-scale
selection, data splits, preservation of adopted Slurm jobs, reporting
and checkpoint recovery. The existing notebook signal-forwarding test
stalled and was excluded from this completed test run
(\texttt{\seqsplit{-k 'not papermill'}}). Official DDA also passed a
strict 537-key CPU checkpoint load. The real GPU smoke check and
experimental metrics are recorded by the submitted preparation and
experiment jobs.

\subsubsection{Long SD1.5 attacker
experiment}\label{record-docs-research-guide-md-long-sd1.5-attacker-experiment}

The current DDPO run trains SD1.5 UNet LoRA against the warmed DDA
defender for 12,800 generated training images, retaining checkpoints at
6,400 and 12,800. It uses eight images per rollout, historical prompt
statistics after at least 16 observations, and reward weights 0.64
evasion, 0.26 CLIP and 0.10 bounded initial-model reference similarity.
Training and evaluation share the same float32 latent-trajectory DDIM
sampler.

The immutable prompt split contains 800 training, 200 validation and 500
test prompts. Validation every 800 training images selects checkpoints
while requiring at least 95\% of initial mean CLIP alignment. Final
evaluation compares the validation-selected checkpoint with initial
SD1.5 on the held-out test prompts. Community Forensics is unused.

One defender update is queued after final attacker evaluation. It reuses
the existing \texttt{\seqsplit{research.defend()}} workflow: the same
selected attacker generates 1,000 fitting images, with 800 paired with
ArtiFact real images for training and 200 pairs for validation. Warmed
DDA receives three epochs of LoRA/head training at learning rate 2e-5.
Before and after updating, the defender is scored on the exact saved
base, attacked and human test images. Both fixed 0.5 and the warmed
defender's validation threshold are retained. The before/after report is
linked from \href{records/docs/image_defender_update_results.tex}{the
defender update results}.

\begin{Verbatim}
PYTHONPATH=. uv --no-cache run --no-sync python scripts/image_experiments.py attach-defender PARENT_MANIFEST
PYTHONPATH=. uv --no-cache run --no-sync python scripts/image_experiments.py submit DEFENDER_MANIFEST --held
\end{Verbatim}

The standalone defender submission adopts the existing final evaluation
job as its dependency, validates its \texttt{\seqsplit{test.json}}
completion receipt, and leaves the original manifest and frozen source
unchanged. Its generation adapter loads the selected long-run checkpoint
and saved test images; training and evaluation remain in the existing
implementation. If Slurm has purged the completed parent from its active
controller, submission can retain verified successful accounting and the
matching completion receipt instead of creating an invalid native
dependency. The worker still checks the parent receipt before training.

Defender integration passed 43 focused image tests, including shared
workflow reuse, identical before/after image sets, generation-only
sampling, adopted completion receipts, restart contracts and reporting.
The GPU defender update is exercised by the queued job after its
predecessor completes. Worker-entry regression checks also exercise the
adopted receipt through the actual job runner, including failed or
mismatched parent receipts.

\href{records/docs/image_attacker_ablation_summary.tex}{Concise long-run
findings} distinguish raw evasion gains from checkpoints meeting the
CLIP requirement. Refresh with:

\begin{Verbatim}
PYTHONPATH=. uv --no-cache run --no-sync python scripts/report_image_attacker_summary.py MANIFEST
\end{Verbatim}

The run and final evaluation remain on regular H100 QoS with checkpoint
recovery. Automatic pre/post reporting jobs are canceled. Reporting is
refreshed manually; the frozen running source, checkpoint configuration
and dependency receipts are preserved for recovery and final evaluation.
Conclusions remain limited to this seed, defender and prompt
distribution.

An additional evaluation-only job can score a saved long-run checkpoint
on the same 500 held-out prompts and noise seeds, reusing the completed
initial-model baseline and human scores. It performs no attacker or
defender training:

\begin{Verbatim}
PYTHONPATH=. uv --no-cache run --no-sync python scripts/image_experiments.py evaluate-checkpoint PARENT_MANIFEST CHECKPOINT
PYTHONPATH=. uv --no-cache run --no-sync python scripts/image_experiments.py submit EVALUATION_MANIFEST --held
\end{Verbatim}

This job requests one hour on regular H100 QoS and records evasion, CLIP
alignment/retention, AUROC and human false positives. Results are linked
from \href{records/docs/image_final_checkpoint_test_results.tex}{the
additional checkpoint evaluation}. This is a post-hoc test of a
requested checkpoint; it does not replace the original
validation-selected attacker or select further checkpoints using test
scores.

\subsubsection{MAML fast
adaptation}\label{record-docs-research-guide-md-maml-fast-adaptation}

\texttt{\seqsplit{MAMLAdaptation}} learns text attacker LoRA starting
weights that work well after five updates against a detector. Each
iteration samples up to three detectors. Every detector branch starts
from the \textbf{same weights [U+03B8]}, practices on \emph{support} prompts,
then computes a loss on different \emph{query} prompts using that same
detector. For example:

\begin{Verbatim}
flowchart TD
    S["Same starting weights [U+03B8]"] --> B["Binoculars: 5 updates on support prompts"]
    S --> R["RoBERTa: 5 updates on support prompts"]
    S --> F["Fast-DetectGPT: 5 updates on support prompts"]
    B --> BQ["Adapted weights [U+03B8][U+2032]B: query loss against Binoculars"]
    R --> RQ["Adapted weights [U+03B8][U+2032]R: query loss against RoBERTa"]
    F --> FQ["Adapted weights [U+03B8][U+2032]F: query loss against Fast-DetectGPT"]
    BQ --> M["Average the three query losses"]
    RQ --> M
    FQ --> M
    M --> U["Backpropagate through practice updates; update original [U+03B8] once"]
    U --> N["Next iteration starts from updated [U+03B8]"]
\end{Verbatim}

\textbf{Support prompts determine the adaptation; query prompts
determine how to improve its starting point.} Query prompts contribute
to training, not final independent evaluation. The adapted branch
weights are discarded; the shared starting weights are retained and
saved as \texttt{\seqsplit{meta\_init}}.

The CLI evaluates adaptation to a held-out detector against the exact
pre-meta starting weights. By default, second-order MAML differentiates
through the practice gradients; \texttt{\seqsplit{--first-order}} uses
FOMAML to reduce compute and memory.

MAML checks the default detector families and all registered text
notebook classifiers. A detector joins the pool only after producing a
finite AI score on a startup probe.
\texttt{\seqsplit{detector\_availability.json}} records every candidate
and unavailable resource; \texttt{\seqsplit{prompt\_manifest.json}}
records the actual training pool and the held-out detector. Style
Embedding needs labeled support centroids; DivEye uses the authors'
released 83-feature DivEye/BiScope classifier and its matching Falcon-7B
and Gemma-1.1-2B-Instruct extractors. Gemma access must be enabled on
Hugging Face or an authorized local checkpoint supplied.

Frozen MAML detectors run on CPU, reserving the single H100 for the
trainable generator and its second-order graph. The Slurm controller
requests 96 physical CPUs for this stage to provide host memory for the
expanded pool. The original 100 episodes, five inner updates and at most
three tasks per episode remain.

\subsubsection{Repeated attacker and defender
adaptation}\label{record-docs-research-guide-md-repeated-attacker-and-defender-adaptation}

Both sides move, so a single attacker or a single retrain is not the
story.

\texttt{\seqsplit{ArmsRaceExperiment}} runs
\texttt{\seqsplit{num\_rounds}} rounds on a fixed evaluation prompt
split that neither side trains on. Before round 1, the initial generator
produces continuations on training prompts to warm up the defender
against matched human references. Warm-up uses the existing defender
budget (up to 2,000 total examples, 3 epochs by default), with an
internal 80/20 training/validation split and the best validation-loss
checkpoint retained. The attacker is not updated during warm-up. The
warmed-up defender enters the attacker reward ensemble before the first
attacker update. Warm-up checkpoints are saved under
\texttt{\seqsplit{defender/warmup/}} and its internal validation metrics
in \texttt{\seqsplit{defender\_warmup.json}}; race rounds keep their
numbering starting at 1. The race reuses the attacker's held-out prompt
split (\texttt{\seqsplit{TextEvasionConfig.evaluation\_fraction}},
default 0.2).

Warm-up protocol and configuration checks were validated with
\texttt{\seqsplit{PYTHONPATH=. uv run --with pytest python -m pytest ai\_content\_detector/tests/test\_equilibrium\_math.py ai\_content\_detector/tests/test\_config.py -q}}.
Full GPU training was not validated locally. GRPO contract test
collection was blocked locally by PyArrow requiring NumPy 2.0 or newer
while NumPy 1.26.3 was installed.

Within each round it takes three paired measurements on the \emph{same}
prompts with the \emph{same} sampling seed:

{\scriptsize\setlength{\tabcolsep}{3pt}\def\LTcaptype{none} % do not increment counter
\begin{longtable}[]{@{}
  >{\raggedright\arraybackslash}p{(\linewidth - 2\tabcolsep) * \real{0.5000}}
  >{\raggedright\arraybackslash}p{(\linewidth - 2\tabcolsep) * \real{0.5000}}@{}}
\toprule\noalign{}
\begin{minipage}[b]{\linewidth}\raggedright
measurement point
\end{minipage} & \begin{minipage}[b]{\linewidth}\raggedright
what it means
\end{minipage} \\
\midrule\noalign{}
\endhead
\bottomrule\noalign{}
\endlastfoot
before attacker update & the attacker's current evasion against the
current defender \\
after attacker update & how much the attacker gained \\
after defender update & how much of that gain the defender recovered \\
\end{longtable}
}

The attacker step budget is enforced.
\texttt{\seqsplit{AdaptiveClassifierDefender}} then retrains a RoBERTa
classifier on fresh attacker outputs paired with matched human
references, and the retrained detector is pushed back into the
attacker's detector ensemble, which closes the loop. The following
numbers are illustrative, not verified results from a recorded run. Each
value is \texttt{\seqsplit{attack\_success\_rate}}: the fraction of
generated evaluation texts with defender AI probability below 0.5.
Higher values mean more successful evasion.

{\scriptsize\setlength{\tabcolsep}{3pt}\def\LTcaptype{none} % do not increment counter
\begin{longtable}[]{@{}
  >{\raggedright\arraybackslash}p{(\linewidth - 6\tabcolsep) * \real{0.2500}}
  >{\raggedright\arraybackslash}p{(\linewidth - 6\tabcolsep) * \real{0.2500}}
  >{\raggedright\arraybackslash}p{(\linewidth - 6\tabcolsep) * \real{0.2500}}
  >{\raggedright\arraybackslash}p{(\linewidth - 6\tabcolsep) * \real{0.2500}}@{}}
\toprule\noalign{}
\begin{minipage}[b]{\linewidth}\raggedright
round
\end{minipage} & \begin{minipage}[b]{\linewidth}\raggedright
before attacker
\end{minipage} & \begin{minipage}[b]{\linewidth}\raggedright
after attacker
\end{minipage} & \begin{minipage}[b]{\linewidth}\raggedright
after defender
\end{minipage} \\
\midrule\noalign{}
\endhead
\bottomrule\noalign{}
\endlastfoot
1 & 0.18 & 0.71 & 0.49 \\
2 & 0.31 & 0.83 & 0.62 \\
3 & 0.44 & 0.88 & 0.71 \\
\end{longtable}
}

In this example, attacker updates increase evasion, while defender
updates reduce it by 22, 21, and 17 percentage points. These numbers do
not establish convergence or non-convergence.

The experiment records paired repeated-update measurements, not a Nash
equilibrium or a Nash gap. \texttt{\seqsplit{RADARDefender}} is only a
compatibility alias; this is not the RADAR alternating paraphraser and
detector algorithm. The generic image arms-race entry point remains
rejected; the dedicated DDA image driver above implements the specified
single-round SD1.5 protocol separately.

\subsection{Evaluation and leakage
controls}\label{record-docs-research-guide-md-evaluation-and-leakage-controls}

\subsubsection{Data flow}\label{record-docs-research-guide-md-data-flow}

Text and image prompt sets are split deterministically before training.
The evaluation portion is immutable. MAML additionally separates support
prompts, used for inner adaptation, from query prompts, used for its
meta-objective.

Benchmark datasets validate one-to-one alignment between prompts, AI
texts, human references, and sample identifiers. They record a source
revision when the dataset exposes one. Trained GRPO and SPIN checkpoints
record method, optimized detector names, and seed metadata. The
benchmark rejects an attack checkpoint whose optimized detector names
overlap declared held-out detectors. Training sample IDs are not yet
stored in attack checkpoints, so sample-level overlap cannot yet be
rejected automatically.

Methods declare whether they rewrite an existing text, generate a new
text from a prompt, or support both. Incompatible dataset and method
pairs are reported as skipped. Failed transformations are counted and
excluded; the original text is never substituted as a successful evasion
output.

\subsubsection{Metrics}\label{record-docs-research-guide-md-metrics}

\begin{itemize}
\tightlist
\item
  \textbf{AUROC} measures how often a random AI example receives a
  higher score than a random human example across all thresholds. It
  requires both classes.
\item
  \textbf{False-positive rate (FPR)} is the fraction of human examples
  incorrectly labeled AI at a threshold.
\item
  \textbf{True-positive rate (TPR)} is the fraction of AI examples
  correctly labeled AI at that threshold.
\item
  \textbf{TPR at a target FPR} uses the best empirically achievable
  threshold at or below the requested FPR. It does not interpolate an
  unattainable operating point.
\item
  \textbf{Attack success rate} is the fraction of generated examples
  whose mean evasion reward exceeds \texttt{\seqsplit{0.5}},
  equivalently whose mean detector AI score is below
  \texttt{\seqsplit{0.5}}. It does not require that the original example
  was first detected.
\item
  \textbf{Semantic similarity} or \textbf{CLIP alignment} checks that
  successful evasion did not simply discard the requested content.
\end{itemize}

Empty, misaligned, non-finite, and invalid single-class metric inputs
are rejected rather than turned into plausible-looking results.

\subsubsection{Feature and notebook leakage
controls}\label{record-docs-research-guide-md-feature-and-notebook-leakage-controls}

\texttt{\seqsplit{FeatureBank}} selects probes on a group-disjoint
validation partition and fits its scaler on training data only. The
explainability notebook assigns stable document identifiers before
paragraph expansion, uses document-disjoint train, validation, and test
sets, restores sentence segmentation, labels SHAP direction for the
actual class, and reserves test data for final evaluation.

\subsection{Running
experiments}\label{record-docs-research-guide-md-running-experiments}

The unified entry point is:

\begin{Verbatim}
PYTHONPATH=. uv run python -m ai_content_detector.rl_evasion.run_experiments \
  --experiment <name>
\end{Verbatim}

{\scriptsize\setlength{\tabcolsep}{3pt}\def\LTcaptype{none} % do not increment counter
\begin{longtable}[]{@{}
  >{\raggedright\arraybackslash}p{(\linewidth - 4\tabcolsep) * \real{0.3333}}
  >{\raggedright\arraybackslash}p{(\linewidth - 4\tabcolsep) * \real{0.3333}}
  >{\raggedright\arraybackslash}p{(\linewidth - 4\tabcolsep) * \real{0.3333}}@{}}
\toprule\noalign{}
\begin{minipage}[b]{\linewidth}\raggedright
Experiment
\end{minipage} & \begin{minipage}[b]{\linewidth}\raggedright
Example
\end{minipage} & \begin{minipage}[b]{\linewidth}\raggedright
Main requirement and output
\end{minipage} \\
\midrule\noalign{}
\endhead
\bottomrule\noalign{}
\endlastfoot
Smoke test & \texttt{\seqsplit{--experiment smoke\_test}} & CPU wiring
check; no training \\
GRPO text &
\texttt{\seqsplit{--experiment grpo\_text --model Qwen/Qwen3.5-9B-Base --epochs 3}}
& Local causal LM, TRL, detectors, GPU; saves a final adapter and
metadata \\
SPIN & \texttt{\seqsplit{--experiment multispin --epochs 5}} & Local
causal LM, CNN/DailyMail, GPU; saves one checkpoint per iteration \\
DDPO image & \texttt{\seqsplit{--experiment ddpo\_image --epochs 50}} &
Stable Diffusion, image detectors, GPU; saves diffusion LoRA
checkpoints \\
Arms race &
\texttt{\seqsplit{--experiment arms\_race --modality text --rounds 10}}
& Working text attacker and adaptive defender; saves round history \\
Meta-adaptation &
\texttt{\seqsplit{--experiment meta\_adapt --epochs 100}} & At least two
functioning detectors, LoRA model, GPU; saves
\texttt{\seqsplit{meta\_init}} \\
Benchmark &
\texttt{\seqsplit{--experiment benchmark --dataset hc3 --epochs 200}} &
Here \texttt{\seqsplit{--epochs}} is used as the maximum sample count;
saves benchmark JSON \\
\end{longtable}
}

Shared options are \texttt{\seqsplit{--model}},
\texttt{\seqsplit{--epochs}}, \texttt{\seqsplit{--rounds}},
\texttt{\seqsplit{--modality}}, \texttt{\seqsplit{--first-order}},
\texttt{\seqsplit{--dataset}}, and \texttt{\seqsplit{--output-dir}}.
Configuration dataclasses in \texttt{\seqsplit{rl\_evasion/config.py}}
expose the remaining learning rates, batch sizes, reward weights, prompt
splits, seeds, and checkpoint intervals.

For finer benchmark control, including explicit held-out detectors and
checkpoints, call the benchmark module directly:

\begin{Verbatim}
PYTHONPATH=. uv run python \
  -m ai_content_detector.rl_evasion.benchmarking.benchmark \
  --dataset hc3 \
  --max-samples 200 \
  --detectors binoculars fast_detect_gpt \
  --heldout-detectors diveye \
  --grpo-checkpoint results/text_evasion/grpo/final \
  --output-dir results/benchmark
\end{Verbatim}

OpenRouter baselines are added only when
\texttt{\seqsplit{OPENROUTER\_API\_KEY}} is present. API failures are
counted as failures and never fall back to returning the prompt.

\subsection{Required models and external
resources}\label{record-docs-research-guide-md-required-models-and-external-resources}

{\scriptsize\setlength{\tabcolsep}{3pt}\def\LTcaptype{none} % do not increment counter
\begin{longtable}[]{@{}
  >{\raggedright\arraybackslash}p{(\linewidth - 2\tabcolsep) * \real{0.5000}}
  >{\raggedright\arraybackslash}p{(\linewidth - 2\tabcolsep) * \real{0.5000}}@{}}
\toprule\noalign{}
\begin{minipage}[b]{\linewidth}\raggedright
Component
\end{minipage} & \begin{minipage}[b]{\linewidth}\raggedright
What must be supplied
\end{minipage} \\
\midrule\noalign{}
\endhead
\bottomrule\noalign{}
\endlastfoot
ModernBERT and TF-IDF & Artifacts under the existing
\texttt{\seqsplit{ml\_pipeline}} \texttt{\seqsplit{bench\_aitextdetect}}
checkpoint registry \\
EfficientNet, CLIP, and DINOv2 classifiers & Artifacts under
\texttt{\seqsplit{bench\_imai\_artifact}} \\
Binoculars and Fast-DetectGPT & Their Hugging Face causal language
models; CUDA is required by availability checks \\
DivEye & Authors' released XGBoost classifier, Falcon-7B and gated
Gemma-1.1-2B-Instruct; all frozen during inference \\
Disrupt-and-Recover & Recovery callable or OpenAI-compatible API, plus a
BERTScore transformer checkpoint \\
IPAD & Phi-3 base model, three released adapters, and a large GPU \\
LUAR style detector & LUAR weights plus task-specific human and AI
support sets \\
MLEP & Trained ResNet-50 classifier checkpoint \\
WaRPAD & DINOv2 ViT-L/14 weights, plus a validation threshold or
calibrator for classification \\
Denoising trajectory & Stable Diffusion and CLIP ViT-L/14 weights \\
GRPO, SPIN, MAML, and DDPO & Local trainable base models, datasets,
functioning reward detectors, and GPU memory \\
\end{longtable}
}

Useful optional environment variables are:

\begin{Verbatim}
export OPENROUTER_API_KEY=...   # OpenRouter benchmark or D&R recovery
export OPENAI_API_KEY=...       # OpenAI-compatible D&R recovery
export OPENAI_BASE_URL=...      # Alternate OpenAI-compatible endpoint
export DR_RECOVERY_MODEL=...    # Recovery model ID for D&R
\end{Verbatim}

Do not commit keys. Model and dataset downloads also require network
access and, for gated models, the provider's authentication and license
acceptance.

\subsection{Extending the
project}\label{record-docs-research-guide-md-extending-the-project}

\subsubsection{Add a
detector}\label{record-docs-research-guide-md-add-a-detector}

\begin{enumerate}
\def\labelenumi{\arabic{enumi}.}
\tightlist
\item
  Subclass \texttt{\seqsplit{BaseDetector}} in the appropriate detector
  module.
\item
  Set \texttt{\seqsplit{name}} and \texttt{\seqsplit{modality}}.
\item
  Implement \texttt{\seqsplit{detect()}} and return a
  \texttt{\seqsplit{DetectionResult}} with an AI-oriented score.
\item
  Put the native statistic, threshold, class order, and
  \texttt{\seqsplit{score\_semantics}} in \texttt{\seqsplit{details}}.
\item
  Implement a truthful \texttt{\seqsplit{is\_available()}} and rely on
  \texttt{\seqsplit{unload()}} or extend its resource list for heavy
  state.
\item
  Add focused tests for boundary decisions, batches, missing resources,
  label ordering, and score direction.
\item
  Register it in the app or experiment detector zoo only after its
  dependencies and failure behavior are explicit.
\end{enumerate}

\subsubsection{Add an evasion
method}\label{record-docs-research-guide-md-add-an-evasion-method}

Declare whether it rewrites existing content or generates from prompts.
Preserve sample identifiers in training metadata, record every detector
optimized by the method, and separate training data from final
evaluation. A failed example must remain a failed example rather than
being replaced by the unmodified input.

\subsection{Retained experiments on Jean
Zay}\label{record-docs-research-guide-md-retained-experiments-on-jean-zay}

The retained attacker is \texttt{\seqsplit{Qwen/Qwen3.5-9B-Base}},
adapted with LoRA against Gemma, ModernBERT-large, or the mean score of
all seven warmed detectors. The retained final checkpoints are Gemma at
1,213 updates, ModernBERT at 1,487, and the ensemble at 502. Their
step-500 snapshots are also retained. A separate RoBERTa step-500
attacker documents the limits of selecting one apparently successful
checkpoint.

{\scriptsize\setlength{\tabcolsep}{3pt}\def\LTcaptype{none} % do not increment counter
\begin{longtable}[]{@{}
  >{\raggedright\arraybackslash}p{(\linewidth - 6\tabcolsep) * \real{0.2500}}
  >{\raggedleft\arraybackslash}p{(\linewidth - 6\tabcolsep) * \real{0.2500}}
  >{\raggedleft\arraybackslash}p{(\linewidth - 6\tabcolsep) * \real{0.2500}}
  >{\raggedleft\arraybackslash}p{(\linewidth - 6\tabcolsep) * \real{0.2500}}@{}}
\toprule\noalign{}
\begin{minipage}[b]{\linewidth}\raggedright
Attacker target
\end{minipage} & \begin{minipage}[b]{\linewidth}\raggedleft
Saved updates
\end{minipage} & \begin{minipage}[b]{\linewidth}\raggedleft
Evaluation texts
\end{minipage} & \begin{minipage}[b]{\linewidth}\raggedleft
Evasion at threshold 0.5
\end{minipage} \\
\midrule\noalign{}
\endhead
\bottomrule\noalign{}
\endlastfoot
Gemma & 1,213 & 500 & 77.2\% \\
ModernBERT-large & 1,487 & 500 & 96.8\% \\
Seven-detector ensemble & 502 & 500 & 94.8\% \\
RoBERTa, diagnostic subset & 500 & 100 & 98.0\% \\
\end{longtable}
}

The RoBERTa diagnostic trajectory later fell to 93\% at update 550 and
36\% at 600 on the same 100 prompts. Its full 500-text evaluation at 600
gave 39.2\%. Only its step-500 model is retained; later measurements
document the fluctuation. A peak does not establish stable convergence.
The three long runs used different training budgets, so their final
scores are not a matched-budget ranking of defender architectures. The
ensemble evaluation was performed at 502, rather than exactly 500.

The \href{records/docs/research_results.tex}{results report} contains
the retained tables, training curves, RoBERTa trajectory, and their
source CSVs and vector SVGs. Other experiment reports and discarded RL
checkpoints have been removed. Notebook models, model/dataset caches,
warmed detectors and recorded data partitions remain available as
execution prerequisites.

The completed experiment performed \textbf{one defender update only} for
each of the three retained long-run attackers. Each job generates
outputs on the existing 1,000 fitting prompts and retrains its warmed
defender(s) with the existing supervised protocol. The 500 evaluation
prompts remain excluded from training. Post-defense evaluation reuses
exactly the saved attack texts, making the before/after comparison
attributable to the defender update. There is no new attacker training
and no second round in these jobs.

\begin{Verbatim}
PYTHONPATH=. uv --no-cache run --no-sync python scripts/retained_races.py prepare
PYTHONPATH=. uv --no-cache run --no-sync python scripts/retained_races.py submit NEW_MANIFEST
PYTHONPATH=. uv --no-cache run --no-sync python scripts/architecture_races.py release NEW_MANIFEST
\end{Verbatim}

All three completed successfully in one development allocation, with no
requeue:

{\scriptsize\setlength{\tabcolsep}{3pt}\def\LTcaptype{none} % do not increment counter
\begin{longtable}[]{@{}
  >{\raggedright\arraybackslash}p{(\linewidth - 8\tabcolsep) * \real{0.2000}}
  >{\raggedleft\arraybackslash}p{(\linewidth - 8\tabcolsep) * \real{0.2000}}
  >{\raggedleft\arraybackslash}p{(\linewidth - 8\tabcolsep) * \real{0.2000}}
  >{\raggedleft\arraybackslash}p{(\linewidth - 8\tabcolsep) * \real{0.2000}}
  >{\raggedright\arraybackslash}p{(\linewidth - 8\tabcolsep) * \real{0.2000}}@{}}
\toprule\noalign{}
\begin{minipage}[b]{\linewidth}\raggedright
Defender
\end{minipage} & \begin{minipage}[b]{\linewidth}\raggedleft
Evasion before update
\end{minipage} & \begin{minipage}[b]{\linewidth}\raggedleft
Evasion after update
\end{minipage} & \begin{minipage}[b]{\linewidth}\raggedleft
Human false positives after
\end{minipage} & \begin{minipage}[b]{\linewidth}\raggedright
Allocation time
\end{minipage} \\
\midrule\noalign{}
\endhead
\bottomrule\noalign{}
\endlastfoot
Gemma & 77.2\% & 0.2\% & 0.2\% & 47m35s \\
ModernBERT-large & 96.8\% & 1.0\% & 0.4\% & 41m45s \\
Seven-detector ensemble & 94.8\% & 0.4\% & 0.0\% & 1h31m13s \\
\end{longtable}
}

These comparisons reuse the same 500 saved attack texts and 500 human
passages at threshold 0.5. Detection recovers against the fixed
attackers; resistance to another adaptive RL attack remains untested.

The updated defenders are also evaluated on the original 500 reserved
base-Qwen outputs and 500 human passages, without generating or training
again. The results report compares their detection and human
false-positive rates with the original warmed detectors, including every
ensemble member. This separate check measures whether learning the RL
outputs preserves detection of the base generator.

All three base-model checks completed successfully:

{\scriptsize\setlength{\tabcolsep}{3pt}\def\LTcaptype{none} % do not increment counter
\begin{longtable}[]{@{}
  >{\raggedright\arraybackslash}p{(\linewidth - 6\tabcolsep) * \real{0.2500}}
  >{\raggedleft\arraybackslash}p{(\linewidth - 6\tabcolsep) * \real{0.2500}}
  >{\raggedleft\arraybackslash}p{(\linewidth - 6\tabcolsep) * \real{0.2500}}
  >{\raggedleft\arraybackslash}p{(\linewidth - 6\tabcolsep) * \real{0.2500}}@{}}
\toprule\noalign{}
\begin{minipage}[b]{\linewidth}\raggedright
Updated branch
\end{minipage} & \begin{minipage}[b]{\linewidth}\raggedleft
Base detection before update
\end{minipage} & \begin{minipage}[b]{\linewidth}\raggedleft
Base detection after update
\end{minipage} & \begin{minipage}[b]{\linewidth}\raggedleft
Human false positives after
\end{minipage} \\
\midrule\noalign{}
\endhead
\bottomrule\noalign{}
\endlastfoot
Gemma & 100.0\% & 100.0\% & 0.2\% \\
ModernBERT-large & 99.0\% & 90.2\% & 0.4\% \\
Seven-detector ensemble & 100.0\% & 98.4\% & 0.0\% \\
\end{longtable}
}

ModernBERT loses some base-model detection at the fixed threshold.
Inside the ensemble, TF-IDF falls from 88.6\% to 28.2\% base detection,
with 7.0\% human false positives and AUROC 0.6158. Recovery against the
fixed RL attacker therefore does not guarantee retention on base
outputs, even when the ensemble remains strong. The report includes
every member and before/after figures.

The three independent jobs each used one H100 and two-hour development
allocations. They save generation, defender optimizer and phase
progress, then requeue after a checkpointed pause. Shared
\texttt{\seqsplit{metrics.csv}} rows identify the branch, phase and
historical attacker budget. See the
\href{records/docs/defender_comparison.tex}{defender-update protocol}
for exact checkpoints, job identities, splits and continuation behavior.
The report selection is explicit, so surviving prerequisite receipts
cannot reintroduce discarded experiments into the Markdown report.

\subsection{Validation
status}\label{record-docs-research-guide-md-validation-status}

The retained-attacker defender workflow passed \textbf{44 focused tests
in 111.55 seconds} in this checkout. They cover zero new attacker
updates, interrupted phase recovery, saved-text evaluation without
regeneration, reserved-prompt validation, cached metric alignment and
report selection. The retained source checkpoint hashes, all
warmed-defender identities and the cached initial metrics were also
checked against the original completed evaluations before release. All
three GPU jobs subsequently completed with exit code zero. The completed
defender checkpoints, identical evaluation-row alignment, and 36 unique
shared metric rows were checked against the saved outputs.

The additional base-model evaluation and its shared submission/reporting
paths passed 26 focused tests in 417.35 seconds. These check use of
saved base outputs and updated weights, completed-result reuse,
rejection of changed inputs or checkpoints, and separate before/after
base-model reporting. The three scoring jobs completed with exit code
zero, no restart, and runtimes of 2m54s, 1m40s and 5m50s for Gemma,
ModernBERT and the ensemble respectively. All 500 paired evaluation
identities and 24 unique shared base-metric rows were verified.
